% !TeX root = ../../surface-lagrangian-torus-fibration.tex
% 2026-11-23 / v1.4: 二領域のsource体積から深度momentの臨界指数1を証明．
\subsection{Canonical potential と actual 球の深部局在}

中央 rescaling 領域そのものの inradius は小さくなり得る．
ここでは固定された physical source tube の中で正の potential を使い，
中央の点を中心とする actual 球が，より広い深部 source 領域から出ないことを示す．
中央 bidisc の内在的な buffer を仮定するものではない．

固定 source 正規形 $xy=tz$ において
\[
 U_t(r,b)=\{|x|<r,\ |y|<r,\ |z|<b,\ xy=tz\}
\]
と置く．以下の半径は全て $t$ に依存せず，必要な source chart の中に
正の余裕を持って含まれるように選ぶ．
$t=\tau^2$ とし，chart 全体の非零正則 $2$-form を
\[
 \Omega_t=\frac{dx\wedge dy}{t}=dX\wedge dY,
 \qquad (X,Y)=(x/\tau,y/\tau)
\]
に固定する．$x\ne0$ では $\Omega_t=d\log x\wedge dz$ であり，
$y\ne0$ では $\Omega_t=-d\log y\wedge dz$ である．
従ってこの frame は二つの座標軸を含めて一つの canonical potential を定める．

\begin{lemma}[固定 source tube 上の正の potential]
\label{ke:source-tube-potential}
十分小さい固定 $r_+,b_+>0$ に対し，その少し外側の固定 chart を用意する．
二つの active plane labels の caps を $\alpha_{x,t},\alpha_{y,t}$ と書くと，
$U_t(r_+,b_+)$ 上で
\begin{equation}\label{ke:eq:source-tube-density}
 c\,\frac{L_t^2}{\alpha_{x,t}^2\alpha_{y,t}^2}|\Omega_t|^2
 \le\omega_t^2\le
 C\,\frac{L_t^2}{\alpha_{x,t}^2\alpha_{y,t}^2}|\Omega_t|^2
\end{equation}
が成り立つ．また $g_t\ge c\,g_{\mathrm{FS}}|_{X_t}$ である．
ある固定 $C_0$ を選べば
\begin{equation}\label{ke:eq:source-positive-potential}
 u_t=\log\frac{\omega_t^2}{|\Omega_t|^2}+2R_t+C_0\ge1,
 \qquad dd^c u_t=\omega_t,
 \qquad \Delta_{g_t}u_t=c_4>0
\end{equation}
となる．最後の $c_4$ は実 Laplacian と K\"ahler 計量の規約だけで決まる定数である．
さらに，tube 内の固定 rescaled bidisc $B_t(c)$ 上で $u_t\le C_c$ である．
\end{lemma}

\begin{proof}
固定 chart では inactive caps と adjunction residue の unit は上下に一様有界である．
関連 KE 主稿の capped density の定義と
\texttt{coarse:cor:degree-d-volume-lower} により
\[
 \Omega_{t,d}^{\mathrm{low}}\asymp
 \frac{L_t^2}{\alpha_{x,t}^2\alpha_{y,t}^2}|\Omega_t|^2,
 \qquad \omega_t^2\ge c\Omega_{t,d}^{\mathrm{low}}.
\]
ここでは元の二つの plane representatives を保つ exact source 座標を用いている．
Cap の定義から，$s_x=-\log|x|^2$，$s_y=-\log|y|^2$ として
\begin{equation}\label{ke:eq:source-cap-depth}
 \alpha_{x,t}\asymp A_0+\min\{s_x,L_t/2\},\qquad
 \alpha_{y,t}\asymp A_0+\min\{s_y,L_t/2\},
 \qquad \alpha_{x,t},\alpha_{y,t}\le C L_t
\end{equation}
である．軸上では対応する $s$ を $+\infty$ と解釈する．
固定 Hermitian frame による有界な深度のずれは比較定数に含める．

上界のため，固定 rescaled 中央 bidisc と，二つの正しい oriented half-arms
によって $U_t(r_+,b_+)$ を覆う．中央部分では命題\ref{ke:source-central-density}，
half-arm では関連 KE 主稿の
\texttt{coarse:lem:pair-arm-density} と
\texttt{coarse:lem:complete-volume-upper} を用いる．
外側 normal 半径と residual 半径を先に大きく取り，
これらの評価をその内部の固定 nested domains 上で適用する．
Half-arm の比較定数は中央 gate まで一様であり，中央では両 caps が
$L_t$ と比較可能なので，三領域の評価は
\eqref{ke:eq:source-tube-density}に統合される．
この被覆には誤った向きの arm の tensor 比較を用いない．

従って $\log(\omega_t^2/\Omega_{t,d}^{\mathrm{low}})$ は tube 全体で上下有界であり，
補題\ref{ke:capped-projective-lower}から射影計量の下界を得る．
\eqref{ke:eq:source-tube-density}の下界と
$\alpha_{x,t},\alpha_{y,t}\le CL_t$ により
$\log(\omega_t^2/|\Omega_t|^2)+2R_t\ge-C$ である．
これで $C_0$ を一度固定して $u_t\ge1$ とできる．
$\Omega_t$ は非零正則 frame なので負 KE 方程式から $dd^c u_t=\omega_t$，
その trace から正定数の Poisson 方程式を得る．
中央 bidisc での上界は命題\ref{ke:source-central-density}そのものである．
\end{proof}

\begin{corollary}[Actual 負 KE 体積に対する capped 前因子の最適性]
\label{ke:source-capped-actual-prefactor-sharp}
固定次数 $d\ge5$ の一般位置全平面退化と
\cite[\texttt{coarse:cor:degree-d-volume-lower},
 \texttt{coarse:prop:capped-prefactor-sharp}]{KELocalCurrent}
の capped functions $\alpha_{i,t}$，canonical density
$\Omega_{0,t}^{(d)}$ を同じものとする．
\[
 \Omega_{t,d}^{(\beta)}
 =L_t^\beta\prod_{i=1}^d\alpha_{i,t}^{-2}\Omega_{0,t}^{(d)}
\]
と置くと，$t$ に一様な正定数 $c$ によって
\[
 \omega_t^2\ge c\,\Omega_{t,d}^{(\beta)}
 \quad\text{on }X_t
\]
が成り立つための必要十分条件は $\beta\le2$ である．
実際 $\beta>2$ なら，任意の単純 source の中央点 $p_t$ に対して
\[
 \frac{\omega_t^2}{\Omega_{t,d}^{(\beta)}}(p_t)
 \asymp L_t^{2-\beta}\longrightarrow0.
\]
この結論は同じ capped family に対する体積下界の前因子の最適性であり，
別の障壁や別の family の最適性を主張しない．
\end{corollary}
\begin{proof}
$\beta\le2$ なら $L_t>1$ で
$\Omega_{t,d}^{(\beta)}=L_t^{\beta-2}\Omega_{t,d}^{\mathrm{low}}
\le\Omega_{t,d}^{\mathrm{low}}$ だから，
上記の大域 capped 下界を適用する．
逆に単純 source $xy=tz$ で
$t=\tau^2$，$(x,y,z)=(\tau X,\tau Y,XY)$ とし，
$p_t$ を $X=Y=0$ の点とする．
命題\ref{ke:source-central-density}から
\[
 \omega_t^2(p_t)\asymp L_t^{-2}|dX\wedge dY|^2.
\]
ここで二つの active plane caps は
$\alpha_x=\alpha_y=A_0+L_t/2$，
残る caps は正の固定定数の間にある．
Adjunction residue の係数は同じ $(X,Y)$ frame で
零から離れた正の上下界を持つので
$\Omega_{t,d}^{(\beta)}(p_t)\asymp
L_t^{\beta-4}|dX\wedge dY|^2$ である．
比を取れば $\beta>2$ の場合の必要条件が従う．
この部分には全域 GH 定理や source の全点距離評価を用いない．
\end{proof}

\begin{proposition}[Source 深部の二つの鋭い体積領域]
\label{ke:source-uniform-deep-volume-tail}
補題\ref{ke:source-tube-potential}の固定 source 正規形 $xy=tz$ で，
$0<r<r_+$，$0<b<b_+<1$ を固定し，
\[
 U_t(r,b)=\{|x|<r,\ |y|<r,\ |z|<b,\ xy=tz\},\qquad
 s_x=-\log|x|^2,\quad s_y=-\log|y|^2
\]
とする（零の座標では深度を $+\infty$ と置く）．
ある固定 $S_0,C>0$ と $t_0>0$ が存在して，全ての
$0<|t|<t_0$ と $S\ge S_0$ に対し
\begin{equation}\label{ke:eq:source-uniform-deep-volume-tail}
 \int_{U_t(r,b)\cap\{\min(s_x,s_y)\ge S\}}
       \frac{\omega_t^2}{2}\le\frac C S
\end{equation}
となる．特に $S=e^A\ge S_0$ と置けば右辺は $Ce^{-A}$ である．
さらに任意の固定 $0<\theta<1/4$ に対して十分大きい $S_0$ を取り直せば，
全ての $S_0\le S\le\theta L$ について
\begin{equation}\label{ke:eq:source-arm-volume-sharp}
 \frac{c_\theta}{S}\le
 \int_{U_t(r,b)\cap\{\min(s_x,s_y)\ge S\}}
       \frac{\omega_t^2}{2}
 \le\frac{C_\theta}{S}
\end{equation}
となる．また固定した $b$ に応じた有限定数 $A_b>0$ を取れば，
全ての $S\ge L/2+A_b$ について
\begin{equation}\label{ke:eq:source-central-exponential-volume}
 c_b L^{-2}e^{-(2S-L)}\le
 \int_{U_t(r,b)\cap\{\min(s_x,s_y)\ge S\}}
       \frac{\omega_t^2}{2}
 \le C_b L^{-2}e^{-(2S-L)}
\end{equation}
となる．定数は固定した source chart，族，次数，半径，
また第一式では $\theta$ に依存してよいが，$t,S$ に依存しない．
両不等式は所定の範囲における体積の次数の二側比較であり，
比が $1$ に収束するという意味の漸近同値ではない．
\end{proposition}

\begin{proof}
$\{z=0\}$ は各滑らかな $X_t$ 内の複素曲線の和であり，
$\omega_t^2$ に関して零測度なので除いて積分する．
$z\ne0$ では $x,y\ne0$ であり，
\[
 s_x+s_y=L+q,\qquad
 L=-\log|t|^2,\quad q=-\log|z|^2>0,\qquad
 \Omega_t=d\log x\wedge dz=-d\log y\wedge dz.
\]
$s_x\le s_y$ とその対称な領域に分ける．前者では
$S\le s_x\le(L+q)/2$ である．
既存の密度上界\eqref{ke:eq:source-tube-density}と
cap 比較\eqref{ke:eq:source-cap-depth}を用いる．
$s_x\le L/2$ なら $s_y\ge L/2$ だから
$\alpha_y\asymp L$，$\alpha_x\asymp A_0+s_x$ であり，
\[
 \omega_t^2\le C(A_0+s_x)^{-2}|d\log x\wedge dz|^2.
\]
$s_x\ge L/2$ なら $s_y\ge s_x$ なので
$\alpha_x,\alpha_y\asymp L$ であり，
\[
 \omega_t^2\le CL^{-2}|d\log x\wedge dz|^2.
\]
後者の $s_x$ 区間の幅は
$\bigl((L+q)/2-\max(S,L/2)\bigr)_+$ 以下である．
$|d\log x\wedge dz|^2$ の積分は，$s_x$，normal angle，
および $z$ の Euclidean area に関する定数倍の積分となる．

$S\le L/2$ の場合，一段目の寄与は
$C\int_S^\infty(A_0+s)^{-2}ds\le C/S$，
二段目は
\[
 CL^{-2}\int_{|z|<b}q(z)\,dA_z\le CL^{-2}\le C/S
\]
である．$S>L/2$ では一段目はなく，
$q\ge Q:=2S-L\ge0$ が必要である．極座標
$dA_z=\pi e^{-q}\,dq$ を使えば二段目は
\[
 CL^{-2}\int_{\substack{q\ge Q\\|z|<b}}(q-Q)e^{-q}\,dq
 \le CL^{-2}e^{-Q}\le C/S.
\]
最後の比較は $S\ge L/2$，$L$ が十分大きいとき
$S e^{-(2S-L)}\le L/2$ による．
対称な $s_y\le s_x$ の部分でも同じ上界が成立する．
固定半径の制約で積分領域を小さくすることは上界を損なわない．
従って全体でも $C/S$ を得る．

次に任意の固定 $0<\theta<1/4$ に対し，
$0<\delta_1<\delta_2<b$ を固定し，$S_0$ を大きく取る．
$S_0\le S\le\theta L$ において
\[
 \delta_1\le|z|\le\delta_2,\qquad S\le s_x\le2S
\]
という領域を取ると，$s_y=L+q-s_x\ge L-2\theta L>L/2$ である．
したがって両normal座標は固定 source tube の内部にあり，
cap は $\alpha_y\asymp L$，$\alpha_x\asymp A_0+s_x$ を満たす．
今度は密度の\emph{下界}を用いると，
\[
 \omega_t^2\ge c(A_0+s_x)^{-2}|d\log x\wedge dz|^2.
\]
固定残余annulus・全角度を積分すると，
\[
 \int_{U_t(r,b)\cap\{\min(s_x,s_y)\ge S\}}
 \frac{\omega_t^2}{2}
 \ge c\int_S^{2S}(A_0+s)^{-2}\,ds\ge\frac{c_\theta}{S}.
\]
先に証明した上界と合わせて
\eqref{ke:eq:source-arm-volume-sharp}を得る．
これは $S\le\theta L$ の全範囲に対し $t,S$ 一様である．

最後に $S\ge L/2+A_b$ とし，$Q=2S-L$ と置く．
固定残余半径について $q_b=-\log b^2$ とし，
$2A_b>q_b+2$ と取れば，$q\ge Q$ の部分は $|z|<b$ に含まれる．
両normal深度が $S\ge L/2$ 以上なので
$\alpha_x,\alpha_y\asymp L$ である．
前半の上界をより正確に書けば，
\[
 \int_{U_t(r,b)\cap\{\min(s_x,s_y)\ge S\}}
 \frac{\omega_t^2}{2}
 \le CL^{-2}\int_Q^\infty(q-Q)e^{-q}\,dq
 =CL^{-2}e^{-Q}.
\]
下界には
\[
 Q+1\le q\le Q+2,\qquad S\le s_x\le S+\tfrac12
\]
の直積領域を用いる．この領域では
$s_y=L+q-s_x\ge S+\tfrac12$ であり，
$|x|,|y|<r$ は小さい $t$ に対して自動的に満たされる．
また $q>q_b$ なので $|z|<b$ である．
密度下界 $\omega_t^2\ge cL^{-2}|d\log x\wedge dz|^2$ を
全normal角度と上記の $s_x,q$ 区間で積分すれば
\[
 \int_{U_t(r,b)\cap\{\min(s_x,s_y)\ge S\}}
 \frac{\omega_t^2}{2}
 \ge cL^{-2}\int_{Q+1}^{Q+2}e^{-q}\,dq
 \ge c_bL^{-2}e^{-Q}.
\]
以上で\eqref{ke:eq:source-central-exponential-volume}が従う．
本証明は固定 source tube の密度二側評価だけを使い，
arm tensor・Gaussian tail・全域 GH 実現や
固定 compact 上の $t\to0$ 収束に依存しない．
\end{proof}

\begin{corollary}[Source 深度 moment の臨界指数]
\label{ke:source-depth-critical-moments}
命題\ref{ke:source-uniform-deep-volume-tail}の固定した
$U_t(r,b)$ と十分大きい固定 $S_0>1$ を取る．
$m_t=\min(s_x,s_y)$ とし，$\gamma>0$ を固定する．
\[
 \mathscr M_{\gamma,t}:=
 \int_{U_t(r,b)\cap\{m_t\ge S_0\}}
   m_t^\gamma\,\frac{\omega_t^2}{2}.
\]
このとき $t$ に依存しない $0<c_\gamma<C_\gamma$ が存在して
\begin{equation}\label{ke:eq:source-depth-critical-moments}
 c_\gamma \mathcal F_\gamma(L_t)
 \le\mathscr M_{\gamma,t}
 \le C_\gamma\mathcal F_\gamma(L_t),\qquad
 \mathcal F_\gamma(L)=
 \begin{cases}
  1,&0<\gamma<1,\\
  \log L,&\gamma=1,\\
  L^{\gamma-1},&\gamma>1
 \end{cases}
\end{equation}
が十分小さい全ての $t\ne0$ で成立する．
従って $T_t=\log m_t$ の source 深部に関する重み付き積分
$\int e^{\gamma T_t}dV_{g_t}$ は，
固定した $\gamma<1$ では一様有界，$\gamma=1$ では
$\log L_t$ の次数で発散する．
この臨界指数 $1$ は既存の密度二側比較からの帰結であり，
全曲面の積分や距離に関する確率法則ではない．
\end{corollary}
\begin{proof}
$V_t(S)$ を命題\ref{ke:source-uniform-deep-volume-tail}の
深部体積とする．$z=0$ の軸は体積零であり，
中央指数評価から任意固定 $\gamma>0$ の積分は有限である．
層別積分の恒等式
\begin{equation}\label{ke:eq:source-moment-layer-cake}
 \mathscr M_{\gamma,t}
 =S_0^\gamma V_t(S_0)
  +\gamma\int_{S_0}^{\infty}S^{\gamma-1}V_t(S)\,dS
\end{equation}
を用いる．任意の固定 $0<\theta<1/4$ を選び，
$L_t$ を十分大きくして $S_0<\theta L_t$ とする．
命題の二側評価から
\[
 V_t(S_0)\asymp1,\qquad
 \int_{S_0}^{\theta L_t}S^{\gamma-1}V_t(S)\,dS
 \asymp\int_{S_0}^{\theta L_t}S^{\gamma-2}\,dS .
\]
後者の積分は $\gamma<1$ で上下から正の固定定数で
抑えられ，$\gamma=1$ では $\log L_t$，
$\gamma>1$ では $L_t^{\gamma-1}$ の次数となる．
これが所要下界を与える．

上界には全 $S\ge S_0$ での $V_t(S)\le C/S$ を
$S\le L_t/2+A_b$ まで用いる．$S\ge L_t/2+A_b$ では
中央指数評価
$V_t(S)\le C_b L_t^{-2}e^{-(2S-L_t)}$ により
\[
 \int_{L_t/2+A_b}^\infty
   S^{\gamma-1}V_t(S)\,dS
 \le C_\gamma L_t^{\gamma-3}
\]
となる．実際 $S=L_t/2+u$ と置き，
指数的に減衰する $u$ について
$(L_t/2+u)^{\gamma-1}$ を積分すればよい．
従って上界も
$1,\log L_t,L_t^{\gamma-1}$ の各次数で成立する．
固定した指数とsource tube以外のパラメータに一様性は要求しない．
\end{proof}

\begin{lemma}[正の Poisson 解の局所対数勾配]
\label{ke:positive-poisson-log-gradient}
実次元 $n$ の Riemannian 多様体で，閉包が定義域に compact に含まれる actual 球
$B_g(p,2\rho)$ 上に
\[
 \operatorname{Ric}(g)\ge-(n-1)K g,\qquad
 u\ge1,\qquad \Delta_g u=c_0>0
\]
が成り立つとする．このとき
\begin{equation}\label{ke:eq:positive-poisson-log-gradient}
 \sup_{B_g(p,\rho)}|\nabla\log u|_g^2
 \le C_n\bigl(\rho^{-2}+K+c_0\bigr).
\end{equation}
球の体積下界，injectivity radius，および $u$ の上界を仮定しない．
\end{lemma}

\begin{proof}
$f=\log u$，$Q=|\nabla f|^2$，$h=c_0/u$ と置くと
\[
 0<h\le c_0,\qquad \Delta f=h-Q,\qquad \nabla h=-h\nabla f.
\]
Bochner の公式と $|\nabla^2 f|^2\ge(\Delta f)^2/n$ により
\begin{align*}
 \Delta Q
 &\ge \frac2n(h-Q)^2
      -2\bigl(h+(n-1)K\bigr)Q
      -2\langle\nabla f,\nabla Q\rangle\\
 &\ge \frac1n Q^2
      -C_n(c_0+K)Q-C_nc_0^2
      -2\langle\nabla f,\nabla Q\rangle.
\end{align*}
Distance cutoff $0\le\chi\le1$ を $B(p,\rho)$ 上で $1$，
$B(p,2\rho)$ の境界近くで $0$ とし，
\[
 \frac{|\nabla\chi|^2}{\chi}\le C_n\rho^{-2},\qquad
 \Delta\chi\ge-C_n(\rho^{-2}+K)
\]
と取る．これは Ricci 下界の Laplacian comparison による．
Cut locus では距離の barrier，または通常の Calabi の小移動を用いる．
$G=\chi Q$ の正の最大点で $\nabla Q=-Q\nabla\chi/\chi$ を代入し，
$\Delta G\le0$ にさらに $\chi$ を掛けると
\[
 0\ge \frac1n G^2-C_n\rho^{-1}G^{3/2}
       -C_n(c_0+K+\rho^{-2})G-C_nc_0^2.
\]
Young の不等式で $G^{3/2}/\rho$ を $G^2/(2n)+C_n\rho^{-4}$ に吸収すると
$G\le C_n(\rho^{-2}+K+c_0)$ を得る．$\chi=1$ の内球で結論となる．
\end{proof}

\begin{proposition}[中央 source の actual 球は深部に留まる]
\label{ke:source-potential-buffer}
固定した $c_1$ と
\[
 c_1^2<b<\beta
\]
を取り，これらを含む固定 source tube とさらに外側の tube を用意する．
$\beta$ は命題\ref{ke:source-shell-gaussian-tail}の shell の半径でよい．
ある $\delta,\kappa>0$ が存在し，十分小さい全ての $t$ と
全ての $p\in\overline{B_t(c_1)}$ に対して
\begin{equation}\label{ke:eq:source-potential-buffer}
 B_{g_t}(p,\delta)
 \subset\mathcal D_{t,\kappa}(b)
 :=\{s_x\ge\kappa L_t,\ s_y\ge\kappa L_t,\ |z|<b\}
\end{equation}
が成り立つ．右辺は上で固定した source chart 内の集合である．
球は全曲面の ambient 距離で定義し，軸上の深度は $+\infty$ とする．
定数は固定半径，族と次数に依存してよいが，$p,t$ に依存しない．
\end{proposition}

\begin{proof}
正の固定幅を持つ三つの nested physical tubes を取り，最内側の residual 半径を
$b$ とし，その中に $\overline{B_t(c_1)}$ を入れる．
固定 ambient 座標 $x,y,z$ の微分は Fubini--Study 計量について有界であり，
外側 tube 上で前補題の $g_t\ge c g_{\mathrm{FS}}$ が成り立つ．
従って，ある $\rho>0$ に対し，中間 tube の各点の $B_{g_t}(q,2\rho)$ は
外側 tube に留まる．実際，最初の退出までに $x,y,z$ の一つが固定正幅だけ
変化するため，その部分だけで固定正の長さを要する．
同じ理由により $\delta>0$ を十分小さく取ると，中央の各 $p$ から長さ
$2\delta$ 未満の曲線は最内側 tube に留まり，特に $|z|<b$ を保つ．

外側 tube の $u_t\ge1$，$\Delta u_t=c_4$ と
$\operatorname{Ric}(g_t)=-g_t$ に補題\ref{ke:positive-poisson-log-gradient}を用いると，
中間 tube 上で $|\nabla\log u_t|\le C$ である．
中央では $u_t(p)\le C$ なので，$q\in B_{g_t}(p,\delta)$ へ至る長さ
$2\delta$ 未満の経路に沿って積分して
\[
 u_t(q)\le u_t(p)e^{2C\delta}\le C'
\]
を得る．経路は上の first-exit 評価により勾配評価の領域内にある．

一方，\eqref{ke:eq:source-tube-density}の下界から
\[
 u_t(q)\ge
 2\log\frac{L_t^2}{\alpha_{x,t}(q)\alpha_{y,t}(q)}-C.
\]
従って $\alpha_{x,t}(q)\alpha_{y,t}(q)\ge cL_t^2$．
各 cap が $CL_t$ 以下なので，両方とも $c'L_t$ 以上である．
\eqref{ke:eq:source-cap-depth}を用い，必要なら $\kappa$ を小さくすると
$s_x(q),s_y(q)\ge\kappa L_t$ を得る．これで全 actual 球について包含が従う．
\end{proof}
